%=============================================================================
% AGENT 8: n_eff IS GRAVITY --- THE IDENTITY THESIS
% The psi-bubble does not perturb gravity; it IS gravity.
% Radiation-reaction analysis for rotating EM modes.
%=============================================================================

\documentclass[11pt]{article}
\usepackage{amsmath,amssymb,physics,geometry,booktabs,tcolorbox}
\usepackage{tikz}
\geometry{margin=1in}

\newcommand{\psifield}{\psi}
\newcommand{\psieff}{\psi_{\rm eff}}
\newcommand{\neff}{n_{\rm eff}}
\newcommand{\etac}{\eta_c}
\newcommand{\avec}{\mathbf{a}}
\newcommand{\kalpha}{k_\alpha}

\title{The $\neff$--Gravity Identity:\\
The $\psi$-Bubble Does Not Perturb Gravity---It \emph{Is} Gravity\\[6pt]
\large Radiation-Reaction Analysis for Rotating EM Modes}
\author{Agent 8 --- DFD $\psi$-Bubble Drive Iteration}
\date{March 2026}

\begin{document}
\maketitle

\begin{abstract}
We develop the thesis that in DFD, the $\psi$-bubble created by
above-threshold EM energy does not \emph{perturb} the gravitational
field---it \emph{is} the gravitational field. The refractive index
$n = e^\psi$ and the gravitational potential $\Phi = -c^2\psi/2$ are not
separate quantities with one ``sourcing'' the other; they are the
\emph{same} physical degree of freedom written in different notation.
This identity, traced directly to the DFD formalism (Section~2 of the
Unified Theory), has three consequences for the $\psi$-bubble drive:
(i)~the static self-force cancellation (Newton's third law) is exact and
inescapable; (ii)~a \emph{dynamic} asymmetric $\psi$-bubble with
time-varying multipole moments radiates $\psi$-waves carrying momentum,
producing a radiation-reaction force analogous to the Abraham--Lorentz
force in electrodynamics; (iii)~this radiation-reaction force is
suppressed by $(v_{\rm rot}/c)^5$ relative to the static gradient force,
rendering it negligible for laboratory devices but potentially significant
for ultra-relativistic configurations. The analysis identifies the
precise conditions under which a rotating EM mode produces net thrust
through the radiation-reaction channel, and quantifies the force for
representative device parameters.
\end{abstract}

\tableofcontents

%=============================================================================
\section{The Identity Thesis}\label{sec:identity}
%=============================================================================

\subsection{The Conceptual Error in Previous Analyses}

Previous treatments of the $\psi$-bubble drive have implicitly adopted a
two-layer picture:
\begin{enumerate}
\item There is a ``gravitational field'' (described by $\psi$, sourced by mass).
\item The EM energy ``perturbs'' this field, creating a $\psi$-bubble
  ``on top of'' gravity.
\item Forces arise from the interaction between the perturbation and the
  background.
\end{enumerate}

This picture, while natural from a GR-trained intuition, is incompatible
with the DFD formalism. In DFD, there is no separate ``gravitational
field'' for the $\psi$-bubble to perturb. The refractive index $n = e^\psi$
\emph{is} the complete description of gravity. When EM energy modifies
$\psi_{\rm eff}$, it is not perturbing gravity---it is \emph{being}
gravity.

\subsection{The Formal Identity}

From Section~2 of the Unified Theory, the DFD formalism establishes:

\paragraph{The optical metric.}
\begin{equation}\label{eq:optical-metric}
d\tilde{s}^2 = -\frac{c^2\,dt^2}{n^2} + d\mathbf{x}^2, \qquad
n(\mathbf{x},t) = e^{\psifield(\mathbf{x},t)}.
\end{equation}

\paragraph{The matter coupling.}
The physical metric governing matter trajectories is:
\begin{equation}\label{eq:physical-metric}
\tilde{g}_{\mu\nu} = \text{diag}\!\left(-c^2 e^{-\psi},\;
e^{+\psi},\; e^{+\psi},\; e^{+\psi}\right),
\end{equation}
from which the non-relativistic equation of motion follows:
\begin{equation}\label{eq:eom}
\frac{d^2\mathbf{x}}{dt^2} = \frac{c^2}{2}\nabla\psi = \avec.
\end{equation}

\paragraph{The Newtonian potential.}
The effective Newtonian potential is:
\begin{equation}\label{eq:potential}
\Phi = -\frac{c^2}{2}\psi.
\end{equation}

\paragraph{The identity.}
Equations~\eqref{eq:optical-metric}--\eqref{eq:potential} establish a
chain of strict equalities:
\begin{equation}\label{eq:chain}
\boxed{
\text{refractive index } n \;=\; e^\psi \;=\;
e^{-2\Phi/c^2} \;\longleftrightarrow\; \text{gravitational potential } \Phi
}
\end{equation}

There is \emph{one} degree of freedom. Whether you call it ``refractive
index,'' ``$\psi$-field,'' or ``gravitational potential,'' you are naming
the same physical quantity. The optical metric \emph{is} the gravitational
metric. Fermat's principle \emph{is} the geodesic equation.

\subsection{Consequence for the $\psi$-Bubble}

When above-threshold EM energy modifies the effective field:
\begin{equation}\label{eq:psieff}
\psieff = \psi_{\rm mass} + \kappa(\eta - \etac)\,\Theta(\eta - \etac),
\end{equation}
the second term is not a ``perturbation on top of gravity.'' It is an
additional contribution to gravity, on exactly the same footing as
$\psi_{\rm mass}$. Matter in the gradient of $\psieff$ accelerates with:
\begin{equation}
\avec = \frac{c^2}{2}\nabla\psieff,
\end{equation}
because this \emph{is} the gravitational acceleration. No coupling
constants, no interaction vertices, no perturbation theory needed. The
gradient of $\psieff$ is gravity, period.

%=============================================================================
\section{The Static Self-Force Problem}\label{sec:self-force}
%=============================================================================

\subsection{Statement of the Problem}

The identity thesis makes the self-force problem sharper, not weaker.
Consider a device that creates an asymmetric $\psieff$ profile. The device
consists of:
\begin{itemize}
\item Mass $M_d$ (the physical structure)
\item EM energy $U_{\rm EM}$ (the field configuration)
\item Resulting $\psieff$ profile with gradient $\nabla\psieff$
\end{itemize}

The device mass $M_d$ experiences acceleration $\avec = (c^2/2)\nabla\psieff$
evaluated at the device location. But the device also \emph{sources}
$\psieff$ through its EM energy. Can it accelerate in its own field?

\subsection{The GR Analog: Why Masses Don't Self-Accelerate}

In GR, a spherically symmetric mass does not accelerate in its own
gravitational field. This follows from Birkhoff's theorem: the field
inside a spherical shell vanishes. More generally, the center of mass
of an isolated system moves on a geodesic of the \emph{external}
field---not the total field including self-contributions.

The DFD analog: for a \emph{static, isolated} $\psieff$ configuration,
the total momentum is conserved. The device-plus-field system has zero
initial momentum, so it cannot develop net momentum. This is exact.

\subsection{The Noether Argument}

The DFD action (Section~2, Eq.~(12)):
\begin{equation}
S_{\rm DFD} = S_\psi + S_h + S_{\rm int} + S_{\rm matter}
\end{equation}
is invariant under spatial translations (on the flat Minkowski
background). By Noether's theorem, the total spatial momentum:
\begin{equation}
\mathbf{P}_{\rm total} = \mathbf{P}_{\rm matter} + \mathbf{P}_\psi
+ \mathbf{P}_{\rm EM}
\end{equation}
is conserved. For an initially static system, $\mathbf{P}_{\rm total} = 0$
at all times. No internal rearrangement can change this.

\begin{tcolorbox}[colback=red!5!white, colframe=red!60!black,
  title=Static Self-Force: Exact Cancellation]
\textbf{Result:} A static, asymmetric $\psieff$ configuration produces
\emph{zero} net force on the device that creates it. This is not
approximate---it is a consequence of Noether's theorem applied to the
DFD action.

\textbf{Physical picture:} The device sits at the bottom of its own
$\psieff$ well. The gradient forces on different parts of the device
(mass elements, EM field stresses, interaction terms) sum to exactly
zero.
\end{tcolorbox}

\subsection{Why the ``Different Source'' Argument Fails for Statics}

One might argue: ``The device mass falls in the gradient created by the
device's EM energy. These are different sources, so self-force is
allowed.'' This argument has a subtle flaw.

Consider the energy functional (from Section~2, Eq.~(33)):
\begin{equation}
E[\psi] = \int d^3x\left[\frac{a_*^2}{8\pi G}
W\!\left(\frac{|\nabla\psi|^2}{a_*^2}\right)
+ \frac{c^2}{2}\rho\psi\right].
\end{equation}

The \emph{static} $\psieff$ profile is the one that minimizes $E$
subject to the constraint that the EM energy density has a given
distribution. At the minimum, the gradient of $E$ with respect to
translations vanishes:
\begin{equation}
\frac{\partial E}{\partial X_i} = 0 \quad \text{at equilibrium},
\end{equation}
where $X_i$ are the device center-of-mass coordinates. This is just
the statement that the equilibrium configuration is force-free. The
``different source'' argument fails because both sources (mass and EM)
are included in the energy functional, and the static equilibrium
automatically balances all forces.

\subsection{What This Rules Out}

The static self-force cancellation rules out the simple picture of
Section~1 of the prompt:
\begin{quote}
``A static device creates an asymmetric $\psieff$ bump. Matter (including
the device) accelerates toward the higher-$\psi$ side.''
\end{quote}
This is incorrect for the device itself. External test particles in
the gradient will accelerate (this is the reaction drive of
Section~8.4 of the propulsion paper). But the device, which sources
the gradient, experiences zero net force in the static limit.

%=============================================================================
\section{The Dynamic Escape: Radiation Reaction}\label{sec:radiation-reaction}
%=============================================================================

\subsection{How Dynamics Evade Noether's Constraint}

Noether's theorem constrains the \emph{total} momentum, not the
\emph{partition} between subsystems. A system with zero total momentum
can have:
\begin{equation}
\mathbf{P}_{\rm device} = -\mathbf{P}_{\rm radiation},
\end{equation}
where $\mathbf{P}_{\rm radiation}$ is momentum carried away by
outgoing $\psi$-waves. If the device radiates $\psi$-waves
preferentially in one direction, it recoils in the opposite direction.
Total momentum is conserved, but the device acquires net momentum at
the expense of radiated field momentum.

This is the exact analog of:
\begin{itemize}
\item \textbf{EM radiation reaction:} A charged particle cannot
  self-accelerate in its own Coulomb field, but it \emph{does}
  experience the Abraham--Lorentz radiation-reaction force when it
  radiates.
\item \textbf{GR radiation reaction:} A binary system cannot
  self-accelerate through its own static gravitational field, but it
  \emph{does} inspiral due to gravitational-wave emission carrying
  energy and momentum away.
\end{itemize}

\subsection{The Abraham--Lorentz Analog for $\psi$}

In electrodynamics, the radiation-reaction force on an accelerating
charge is:
\begin{equation}
\mathbf{F}_{\rm AL} = \frac{q^2}{6\pi\epsilon_0 c^3}\dot{\avec}.
\end{equation}
This arises because the radiating charge loses momentum to the EM
field, and by Newton's third law, the field exerts a reaction force
on the charge.

For a $\psi$-field source with time-varying multipole moments, the
analogous radiation-reaction force follows from the momentum flux
in the radiated $\psi$-waves.

\subsection{$\psi$-Wave Radiation from a Time-Varying Source}

The linearized $\psi$ field equation in the high-gradient (Newtonian)
regime is:
\begin{equation}\label{eq:psi-wave}
\Box\psi = \frac{1}{c^2}\ddot{\psi} - \nabla^2\psi
= -\frac{8\pi G}{c^2}\rho_{\rm eff},
\end{equation}
where $\rho_{\rm eff}$ includes both mass density and the effective
EM source from the above-threshold coupling:
\begin{equation}
\rho_{\rm eff} = \rho_{\rm mass}
+ \frac{\kappa}{\alpha/4}\frac{U_{\rm EM}}{c^2}\,\Theta(\eta - \etac).
\end{equation}

The retarded solution is:
\begin{equation}
\psi(\mathbf{x},t) = \frac{2G}{c^2}\int\frac{\rho_{\rm eff}
(\mathbf{x}',t-|\mathbf{x}-\mathbf{x}'|/c)}{|\mathbf{x}-\mathbf{x}'|}
\,d^3x'.
\end{equation}

\subsection{Multipole Expansion of Radiated Power}

For a source with characteristic size $R$ and angular frequency $\omega$,
expanding in powers of $\omega R/c$:

\paragraph{Monopole ($\ell = 0$).}
The monopole moment is the total ``$\psi$-charge'':
\begin{equation}
Q_\psi = \int \rho_{\rm eff}\,d^3x.
\end{equation}
For a device with fixed mass and constant total EM energy,
$\dot{Q}_\psi = 0$. \textbf{No monopole radiation.} (This is the
analog of charge conservation suppressing monopole EM radiation.)

\paragraph{Dipole ($\ell = 1$).}
The dipole moment is:
\begin{equation}
\mathbf{d}_\psi = \int \mathbf{x}\,\rho_{\rm eff}\,d^3x.
\end{equation}
For a system with conserved total momentum, $\ddot{\mathbf{d}}_\psi = 0$
in the center-of-mass frame. \textbf{No dipole radiation.} (This is the
analog of momentum conservation suppressing dipole gravitational-wave
radiation in GR.)

\paragraph{Quadrupole ($\ell = 2$).}
The quadrupole moment is:
\begin{equation}
Q_{ij} = \int \rho_{\rm eff}\!\left(x_i x_j
- \frac{1}{3}\delta_{ij}r^2\right)d^3x.
\end{equation}
Time variation of $Q_{ij}$ \emph{does} radiate. The radiated power is:
\begin{equation}\label{eq:quadrupole-power}
P_\psi = \frac{G}{5c^5}\langle\dddot{Q}_{ij}\dddot{Q}_{ij}\rangle.
\end{equation}

This is the leading radiation channel. For a rotating EM mode, the
quadrupole moment varies at $2\omega$, giving:
\begin{equation}
\dddot{Q}_{ij} \sim (2\omega)^3 Q_0,
\end{equation}
where $Q_0$ is the quadrupole amplitude set by the EM energy
distribution.

%=============================================================================
\section{The Rotating EM Mode: Detailed Calculation}\label{sec:rotating}
%=============================================================================

\subsection{Configuration}

Consider a cylindrical cavity supporting a rotating $\text{TE}_{111}$
mode. The EM energy density rotates with angular frequency $\omega$,
creating a time-dependent above-threshold region that orbits the
cavity axis. The key parameters:
\begin{itemize}
\item Cavity radius: $R$
\item Mode frequency: $\omega$
\item Peak EM energy density: $U_0$
\item Above-threshold volume: $V_+$
\item Ambient matter density: $\rho$
\item Above-threshold ratio: $\eta_0/\etac \equiv \chi$
\end{itemize}

\subsection{The Effective $\psi$ Source}

The rotating mode creates an effective source:
\begin{equation}
\rho_{\rm eff}(\mathbf{x},t) = \rho
+ \frac{\kappa U_0}{c^2}f(\mathbf{x},t)\,\Theta(f - f_c),
\end{equation}
where $f(\mathbf{x},t)$ describes the rotating EM energy profile and
$f_c$ is the threshold value. The quadrupole moment of $\rho_{\rm eff}$
has the form:
\begin{equation}
Q_{ij}(t) = Q_0\begin{pmatrix}
\cos 2\omega t & \sin 2\omega t & 0 \\
\sin 2\omega t & -\cos 2\omega t & 0 \\
0 & 0 & 0
\end{pmatrix},
\end{equation}
with amplitude:
\begin{equation}\label{eq:Q0}
Q_0 \sim \frac{\kappa U_0}{c^2}\,V_+ R^2 \cdot g(\chi),
\end{equation}
where $g(\chi)$ is a geometric factor encoding the shape of the
above-threshold region. For $\chi \gg 1$ (well above threshold),
$g(\chi) \sim 1$.

\subsection{Radiated Power}

From Eq.~\eqref{eq:quadrupole-power}:
\begin{align}
P_\psi &= \frac{G}{5c^5}\langle\dddot{Q}_{ij}\dddot{Q}_{ij}\rangle
\nonumber\\
&= \frac{G}{5c^5}\cdot 2\cdot(2\omega)^6 Q_0^2 \nonumber\\
&= \frac{128\,G\omega^6 Q_0^2}{5c^5}.
\label{eq:Ppsi}
\end{align}

Substituting $Q_0$ from Eq.~\eqref{eq:Q0}:
\begin{equation}\label{eq:Ppsi-expanded}
\boxed{
P_\psi = \frac{128\,G\omega^6}{5c^9}\,\kappa^2 U_0^2
V_+^2 R^4\,g^2(\chi)
}
\end{equation}

\subsection{Momentum Flux and Radiation-Reaction Force}

For an \emph{isotropic} radiator, the net momentum flux vanishes.
A radiation-reaction force requires asymmetric radiation---more
$\psi$-wave power emitted in one direction than the other.

\paragraph{Source of asymmetry.}
The rotating mode is axially symmetric (around the rotation axis) but
can be made asymmetric along the axis through:
\begin{enumerate}
\item Geometric asymmetry of the cavity (open vs.\ shielded ends)
\item Axial gradient in the EM energy profile
\item Asymmetric above-threshold boundary
\end{enumerate}

Define the asymmetry parameter:
\begin{equation}
\mathcal{A} = \frac{P_\psi^+ - P_\psi^-}{P_\psi^+ + P_\psi^-},
\end{equation}
where $P_\psi^\pm$ are the radiated powers in the $+z$ and $-z$
hemispheres. For a device with open/shielded asymmetry
(as in Agent~2's solenoid geometry), $\mathcal{A} \sim (L_o - L_s)/(L_o + L_s)$.

The radiation-reaction force is:
\begin{equation}\label{eq:F-rad}
\boxed{
F_{\rm rad} = \frac{\mathcal{A}\,P_\psi}{c}
= \frac{128\,G\omega^6\,\kappa^2 U_0^2 V_+^2 R^4\,g^2(\chi)
\,\mathcal{A}}{5c^{10}}
}
\end{equation}

%=============================================================================
\section{Numerical Estimates}\label{sec:numerics}
%=============================================================================

\subsection{Laboratory Scale}

Representative parameters for a high-field laboratory device:
\begin{center}
\begin{tabular}{lll}
\toprule
Parameter & Symbol & Value \\
\midrule
Cavity radius & $R$ & 0.1 m \\
Mode frequency & $\omega$ & $2\pi \times 10^9$ Hz \\
Peak EM energy density & $U_0$ & $4 \times 10^7$ J/m$^3$ (10 T) \\
Above-threshold volume & $V_+$ & $10^{-3}$ m$^3$ \\
Self-coupling & $\kappa$ & 51.4 \\
Asymmetry & $\mathcal{A}$ & 0.8 \\
Geometric factor & $g(\chi)$ & 1 \\
\bottomrule
\end{tabular}
\end{center}

Computing $Q_0$:
\begin{align}
Q_0 &\sim \frac{\kappa U_0}{c^2}\,V_+ R^2
= \frac{51.4 \times 4\times10^7}{(3\times10^8)^2}
\times 10^{-3} \times 10^{-2} \nonumber\\
&\approx 2.3 \times 10^{-14}\;\text{kg\,m}^2.
\end{align}

Computing $P_\psi$:
\begin{align}
P_\psi &= \frac{128\,G\omega^6 Q_0^2}{5c^5} \nonumber\\
&= \frac{128 \times 6.67\times10^{-11} \times (6.3\times10^9)^6
\times (2.3\times10^{-14})^2}{5 \times (3\times10^8)^5} \nonumber\\
&= \frac{128 \times 6.67\times10^{-11} \times 6.3\times10^{58}
\times 5.3\times10^{-28}}{5 \times 2.4\times10^{42}} \nonumber\\
&\sim \frac{128 \times 2.2\times10^{20}}{1.2\times10^{43}} \nonumber\\
&\sim 2.4 \times 10^{-21}\;\text{W}.
\end{align}

The radiation-reaction force:
\begin{equation}
F_{\rm rad} = \frac{\mathcal{A}\,P_\psi}{c}
\approx \frac{0.8 \times 2.4\times10^{-21}}{3\times10^8}
\approx 6 \times 10^{-30}\;\text{N}.
\end{equation}

\begin{tcolorbox}[colback=yellow!5!white, colframe=yellow!60!black,
  title=Laboratory Radiation-Reaction Force]
\begin{equation}
F_{\rm rad}^{\rm lab} \sim 10^{-29}\;\text{N}
\end{equation}
This is \textbf{utterly negligible}---roughly $10^5$ times smaller
than the $\psi$-wave recoil computed by Agent~2 ($\sim 10^{-24}$ N),
which was itself identified as undetectable.

\textbf{The suppression factor} is $(v_{\rm rot}/c)^5$ where
$v_{\rm rot} = \omega R \sim 6\times10^8$ m/s (note: this is
superluminal for the quoted parameters, meaning the multipole
expansion breaks down and the actual power is \emph{lower} than
the estimate).
\end{tcolorbox}

\subsection{The $(v/c)^5$ Suppression}

The radiation-reaction force scales as:
\begin{equation}
F_{\rm rad} \propto \frac{G}{c^{10}}\,\omega^6 \propto
\frac{G}{c^{10}}\left(\frac{v_{\rm rot}}{R}\right)^6
= \frac{G}{c^4}\left(\frac{v_{\rm rot}}{c}\right)^6\frac{1}{R^6}.
\end{equation}

The ratio to the static gradient force
$F_{\rm static} \sim (c^2/2)\kappa U_0 V_+ / R$ is:
\begin{equation}
\frac{F_{\rm rad}}{F_{\rm static}} \sim
\frac{G\,\kappa U_0 V_+}{c^5 R^5}\,\omega^6 R^6
\sim \frac{G\,\kappa U_0 V_+}{c^5}\left(\frac{\omega R}{c}\right)^5
\omega.
\end{equation}

For $\omega R / c \lesssim 1$ (non-superluminal rotation), this ratio
contains the factor $(v_{\rm rot}/c)^5 \leq 1$, combined with the
gravitational weakness factor $G U_0 V_+ / c^5$.

This double suppression---gravitational weakness \emph{and} relativistic
velocity ratio---makes radiation reaction negligible for any
laboratory-scale device.

%=============================================================================
\section{The Alcubierre Analogy: Where It Works and Where It Fails}
\label{sec:alcubierre}
%=============================================================================

\subsection{The Analogy}

The prompt draws an analogy with the Alcubierre warp metric: the device
modifies the refractive index around it, and ``falls'' through the
modified medium. In Alcubierre's construction:
\begin{equation}
ds^2 = -c^2 dt^2 + (dx - v_s f(r_s)\,dt)^2 + dy^2 + dz^2,
\end{equation}
the function $f(r_s)$ creates a contraction ahead and expansion
behind, and the enclosed volume ``surfs'' the warp bubble at velocity
$v_s$.

In DFD language, this would correspond to:
\begin{equation}
\psieff(\mathbf{x},t) = \psi_{\rm front}(x - v_s t)
- \psi_{\rm rear}(x - v_s t),
\end{equation}
with $\psi_{\rm front} > \psi_{\rm rear}$, creating a gradient that
pushes the enclosed volume forward.

\subsection{Where the Analogy Fails}

The Alcubierre metric requires \emph{exotic matter} (negative energy
density) to maintain the warp bubble. In DFD, the corresponding
requirement is:

\paragraph{Who maintains the gradient?}
The $\psieff$ gradient is sourced by the EM energy distribution. If the
device moves, the EM energy moves with it (it is bound to the device
hardware). The $\psieff$ profile thus moves with the device, maintaining
the gradient.

\paragraph{But Noether's theorem applies.}
The device-plus-EM-field system has conserved total momentum. For the
device to accelerate, \emph{something} must carry momentum in the opposite
direction. In the Alcubierre case, the exotic matter provides the ``push.''
In the DFD case, the only available momentum carriers are:
\begin{enumerate}
\item Expelled matter (reaction drive---works, per Section~8.4 of the
  propulsion paper)
\item Radiated $\psi$-waves (radiation reaction---works in principle,
  negligible in practice)
\item External $\psi$-field gradient (falling in an external
  gravitational field---not self-propulsion)
\end{enumerate}

There is \emph{no} fourth option. The DFD action is translationally
invariant, momentum is conserved, and the device cannot bootstrap itself
into motion without ejecting matter or radiating field momentum.

\subsection{The Charged-Particle Analogy: Corrected}

The prompt suggests: ``A charged particle doesn't accelerate in its own
Coulomb field, but it DOES experience radiation reaction when it
radiates.'' This is correct---and Section~\ref{sec:radiation-reaction}
develops exactly this analog. But the conclusion is sobering:

\begin{itemize}
\item The EM radiation-reaction force (Abraham--Lorentz) scales as
  $q^2 \dot{a}/(6\pi\epsilon_0 c^3)$. For laboratory charges, this is
  tiny ($\sim 10^{-24}$ N for electron-scale charges with
  $\dot{a} \sim 10^{20}$ m/s$^3$).

\item The $\psi$-radiation-reaction force is \emph{additionally}
  suppressed by $G/c^2$ relative to the EM case, because $\psi$-waves
  couple gravitationally.

\item The ``radiation'' escape route works in principle but is suppressed
  by the same gravitational weakness that makes all
  $\psi$-wave effects negligible at laboratory scale.
\end{itemize}

%=============================================================================
\section{Can the Gradient Be Sustained While Moving?}\label{sec:sustain}
%=============================================================================

\subsection{The Moving $\psieff$ Profile}

Consider a device moving at velocity $v$ through a medium of density
$\rho$. The device carries its EM field configuration, so:
\begin{equation}
\psieff(\mathbf{x},t) = \psi_{\rm mass}(\mathbf{x})
+ \kappa(\eta(\mathbf{x} - \mathbf{v}t) - \etac)
\,\Theta(\eta(\mathbf{x} - \mathbf{v}t) - \etac).
\end{equation}

The $\psieff$ bump moves with the device. Matter in the path of the
device enters the above-threshold region and is accelerated. This is
precisely the reaction drive mechanism.

\subsection{Self-Consistency of the Moving Solution}

For the reaction drive (expelling matter), the moving solution is
self-consistent:
\begin{enumerate}
\item Device creates $\psieff$ gradient in surrounding medium
\item Ambient gas enters the gradient region
\item Gas is accelerated to $v_{\rm exhaust} \sim c\sqrt{\kappa(\eta - \etac)}$
\item Device recoils by momentum conservation
\item Device moves, carrying its $\psieff$ bump to fresh medium
\item Repeat
\end{enumerate}

This is a conventional (if extraordinary) reaction drive. It does
\emph{not} require the Alcubierre mechanism. The device does not
``surf its own wave''---it throws matter backwards and moves forward.

\subsection{The Vacuum Limit}

In perfect vacuum ($\rho \to 0$), the above-threshold condition
$\eta = U_{\rm EM}/(\rho c^2) > \etac$ is trivially satisfied everywhere
the EM field exists. But there is no matter to expel, so the reaction
drive produces zero thrust (infinite $I_{\rm sp}$, zero mass flow rate,
zero force).

In this limit, the \emph{only} thrust mechanism is radiation reaction
(Section~\ref{sec:radiation-reaction}), which is negligible.

This is the fundamental limitation: the $\neff$--gravity identity
guarantees that the $\psieff$ gradient produces real gravitational
acceleration on test masses, but Noether's theorem guarantees that
the device cannot exploit this to self-accelerate without expelling
mass or radiating field momentum.

%=============================================================================
\section{The ``Different Source'' Argument Revisited}\label{sec:different-source}
%=============================================================================

\subsection{Why ``Mass Falls in EM-Sourced Gradient'' Is Incomplete}

The argument: ``The EM energy and the device mass are DIFFERENT sources.
The device mass falls in the $\psieff$ gradient created by the device's
EM energy.''

This is true as a local statement: at any given point, the gravitational
acceleration is indeed $(c^2/2)\nabla\psieff$, and the EM contribution to
$\psieff$ is distinct from the mass contribution. But it is \emph{false}
as a global statement about net force on the device, because:

\paragraph{The EM field is not free.}
The EM energy is stored in the device's cavity/solenoid. The EM field
configuration is maintained by currents in the device hardware. When the
$\psieff$ gradient exerts a force on the device mass pulling it one way,
it simultaneously exerts a force on the EM field energy pulling it the
\emph{other} way (or equivalently, the field stresses on the cavity walls
balance the gravitational pull).

\paragraph{Formal statement.}
Let $\mathbf{F}_{\rm on\,mass}$ be the gravitational force on the device
mass from the EM-sourced $\psieff$, and $\mathbf{F}_{\rm on\,EM}$ be the
reaction force on the EM configuration. By the translational invariance
of the total action:
\begin{equation}
\mathbf{F}_{\rm on\,mass} + \mathbf{F}_{\rm on\,EM}
+ \mathbf{F}_{\rm self\,mass} = 0.
\end{equation}

The device mass is attracted toward the peak of the EM-sourced $\psieff$.
The EM field is attracted toward the peak of the mass-sourced $\psi$.
These forces \emph{compress} the device (internal stress) but do not
produce net translation.

\subsection{The Exception: Radiation}

The one exception, as developed in Section~\ref{sec:radiation-reaction},
is when the system radiates. Radiation carries momentum to infinity,
breaking the balance. But the radiated momentum is:
\begin{equation}
\dot{P}_{\rm rad} = \frac{P_\psi\,\mathcal{A}}{c}
\sim \frac{G}{c^{10}}\,\kappa^2 U_0^2 V_+^2 R^4\,\omega^6\,\mathcal{A},
\end{equation}
which is negligible for laboratory parameters.

%=============================================================================
\section{Summary and Implications}\label{sec:summary}
%=============================================================================

\begin{tcolorbox}[colback=green!5!white, colframe=green!50!black,
  title=Central Results]

\begin{enumerate}
\item \textbf{The Identity:} In DFD, $\neff = e^{\psieff}$ and the
  gravitational potential $\Phi = -c^2\psieff/2$ are the \emph{same}
  degree of freedom. The $\psi$-bubble does not ``perturb gravity''---it
  \emph{is} gravity. This is not an approximation; it is the founding
  postulate of DFD.

\item \textbf{Static self-force cancellation:} Noether's theorem
  (translational invariance of the DFD action) guarantees that a static
  $\psieff$ configuration produces zero net force on its source. This is
  exact and model-independent within DFD.

\item \textbf{Radiation-reaction force exists:} A rotating EM mode
  creates a time-varying quadrupole moment in $\rho_{\rm eff}$, radiating
  $\psi$-waves and producing a radiation-reaction force. This is the
  DFD analog of the Abraham--Lorentz force.

\item \textbf{Radiation reaction is negligible:} The force scales as
  $F_{\rm rad} \sim (G/c^{10})\,\kappa^2 U_0^2 V_+^2 R^4 \omega^6
  \mathcal{A}$, giving $\sim 10^{-29}$ N for laboratory parameters.
  The $(v/c)^5$ suppression plus gravitational weakness makes this
  undetectable.

\item \textbf{The reaction drive remains the viable mechanism:} The
  above-threshold $\psieff$ gradient produces real, large accelerations on
  external matter ($\avec = (c^2/2)\nabla\psieff$ with $\kappa \approx 51$).
  Expelling this accelerated matter gives $I_{\rm sp} \sim 10^7$ s. This
  mechanism is fully consistent with momentum conservation.

\item \textbf{Vacuum self-propulsion is not possible:} In the absence of
  matter to expel, the device can only generate thrust through
  $\psi$-wave radiation reaction, which is negligible. The Alcubierre
  analogy fails because DFD momentum conservation forbids reactionless
  propulsion.
\end{enumerate}
\end{tcolorbox}

\subsection{Implications for Device Design}

\begin{enumerate}
\item \textbf{Optimize for reaction drive, not self-propulsion.}
  The viable mechanism is the extraordinary reaction drive with
  $I_{\rm sp} \sim 10^7$ s. Device design should maximize the
  above-threshold volume, the $\eta/\etac$ ratio, and the mass
  flow rate of propellant through the acceleration region.

\item \textbf{Rotating modes are irrelevant for thrust.}
  The radiation-reaction channel from rotating EM modes is suppressed
  by $(G/c^5)(v_{\rm rot}/c)^5$ relative to the reaction-drive
  mechanism. There is no advantage to rotating the EM mode for
  propulsion purposes.

\item \textbf{The $\neff$--gravity identity is an asset, not a problem.}
  The identity means that the $\psieff$ gradient is \emph{real gravity},
  not a weak perturbation requiring amplification. External matter
  experiences the full $\avec = (c^2/2)\nabla\psieff$ without any
  $G/c^4$ suppression. The bottleneck is getting matter into the
  gradient region and maintaining the above-threshold condition.
\end{enumerate}

\subsection{Open Questions}

\begin{enumerate}
\item \textbf{Nonlinear $\psi$-wave effects:} The quadrupole formula
  used here is the linearized result. For strong $\psieff$ gradients
  ($\kappa(\eta - \etac) \sim 1$), nonlinear $\psi$-wave emission
  could differ qualitatively. This requires numerical integration of
  the full nonlinear field equation.

\item \textbf{Medium-mediated momentum transfer:} In a medium
  (interstellar gas, residual atmosphere), the $\psieff$ gradient
  continuously accelerates ambient particles. If the device moves through
  a medium, it experiences a ``$\psi$-drag'' or ``$\psi$-thrust''
  depending on the geometry. This medium-mediated channel deserves
  separate analysis for interstellar propulsion scenarios.

\item \textbf{Threshold dynamics:} The Heaviside function
  $\Theta(\eta - \etac)$ in Eq.~\eqref{eq:psieff} creates a sharp
  boundary. At this boundary, $\nabla\psieff$ has a discontinuity.
  The physical smoothing of this boundary (over what length scale?
  by what mechanism?) affects the force profile and deserves
  investigation.
\end{enumerate}

\end{document}
